% Seminar 3 - Pagina 2. Fragment pentru
% % Seminar 3 - Pagina 2. Fragment pentru
% % Seminar 3 - Pagina 2. Fragment pentru
% % Seminar 3 - Pagina 2. Fragment pentru
% \input{seminars/pages/seminar3-pagina2-ro}; fără preambul sau pagină nouă.
% NOTE EDITORIALE:
% - Pagina continuă problema 2(b) din Seminarul 3 - Pagina 1.
% - Manuscrisul alternează notațiile P_y/P_z și p_y/p_z;
%   proiecțiile sunt notate consecvent p_Y și p_Z.
% - Diagrama rearanjează schița manuscrisului pentru lizibilitate,
%   păstrând proiecțiile și cele două aplicații componente.
% - O observație intercalată pare să fie: fixez a din Z,
%   i_Y(y)=(y,a). Am păstrat-o, deși nu intervine în demonstrație.
% - Calculele tăiate din dreptul celor două identități sunt omise;
%   pașii finali lizibili sunt păstrați.
% - În câteva locuri manuscrisul scrie identitatea fără a preciza
%   mulțimea pe care acționează; domeniile sunt explicite mai jos.

\subsection*{Seminar 3 - Pagina 2}

\noindent\textbf{(b) Continuare.} Avem corespondența
\[
  \operatorname{Hom}(X,Y\times Z)
  \;\overset{\alpha}{\underset{\beta}{\rightleftarrows}}\;
  \operatorname{Hom}(X,Y)\times\operatorname{Hom}(X,Z).
\]
Fie $p_Y:Y\times Z\to Y$ și $p_Z:Y\times Z\to Z$ proiecțiile
canonice:
\[
 p_Y(y,z)=y,\qquad p_Z(y,z)=z.
\]
Diagrama proiecțiilor și a aplicațiilor componente este
\[
\begin{tikzcd}[column sep=large, row sep=large]
 & X \arrow[dl, "h"'] \arrow[dr, "g"]
     \arrow[d, dashed, "{\beta(g,h)}"] & \\
 Z & Y\times Z \arrow[l, "p_Z"] \arrow[r, "p_Y"'] & Y
\end{tikzcd}
\]
Pentru $f:X\to Y\times Z$, definim
\[
  \alpha(f)=(p_Y\circ f,\,p_Z\circ f).
\]
Reciproc, pentru $g:X\to Y$ și $h:X\to Z$, punem
\[
  \beta(g,h)(x)=(g(x),h(x)),\qquad x\in X.
\]

\medskip
\noindent\emph{Observație intercalată în manuscris.} Fixăm $a\in Z$ și
considerăm
\[
 \iota_Y:Y\longrightarrow Y\times Z,\qquad
 \iota_Y(y)=(y,a).
\]

\medskip
\noindent Arătăm mai întâi că $\alpha\circ\beta=\operatorname{id}$.
Pentru orice $g:X\to Y$ și $h:X\to Z$,
\[
  \alpha\bigl(\beta(g,h)\bigr)=(g,h).
\]
Într-adevăr,
\[
\begin{aligned}
 \bigl(p_Y\circ\beta(g,h)\bigr)(x)
   &=p_Y\bigl(\beta(g,h)(x)\bigr) \\
   &=p_Y(g(x),h(x))=g(x),
 \qquad \forall x\in X,\\[3pt]
 \bigl(p_Z\circ\beta(g,h)\bigr)(x)
   &=p_Z\bigl(\beta(g,h)(x)\bigr) \\
   &=p_Z(g(x),h(x))=h(x),
 \qquad \forall x\in X.
\end{aligned}
\]
Prin urmare,
\[
  (\alpha\circ\beta)(g,h)
    =(p_Y\circ\beta(g,h),\,p_Z\circ\beta(g,h))
    =(g,h).
\]

\medskip
\noindent Reciproc, demonstrăm că $\beta\circ\alpha=\operatorname{id}$.
Fie $f:X\to Y\times Z$. Atunci
\[
 (\beta\circ\alpha)(f)=f.
\]
Pentru $x\in X$, notăm $f(x)=(y,z)$. Avem
\[
\begin{aligned}
 \beta\bigl(\alpha(f)\bigr)(x)
   &=\beta(p_Y\circ f,\,p_Z\circ f)(x)\\
   &=\bigl((p_Y\circ f)(x),\,(p_Z\circ f)(x)\bigr)\\
   &=\bigl(p_Y(y,z),\,p_Z(y,z)\bigr)\\
   &=(y,z)=f(x).
\end{aligned}
\]
Deci cele două corespondențe sunt inverse una alteia.
; fără preambul sau pagină nouă.
% NOTE EDITORIALE:
% - Pagina continuă problema 2(b) din Seminarul 3 - Pagina 1.
% - Manuscrisul alternează notațiile P_y/P_z și p_y/p_z;
%   proiecțiile sunt notate consecvent p_Y și p_Z.
% - Diagrama rearanjează schița manuscrisului pentru lizibilitate,
%   păstrând proiecțiile și cele două aplicații componente.
% - O observație intercalată pare să fie: fixez a din Z,
%   i_Y(y)=(y,a). Am păstrat-o, deși nu intervine în demonstrație.
% - Calculele tăiate din dreptul celor două identități sunt omise;
%   pașii finali lizibili sunt păstrați.
% - În câteva locuri manuscrisul scrie identitatea fără a preciza
%   mulțimea pe care acționează; domeniile sunt explicite mai jos.

\subsection*{Seminar 3 - Pagina 2}

\noindent\textbf{(b) Continuare.} Avem corespondența
\[
  \operatorname{Hom}(X,Y\times Z)
  \;\overset{\alpha}{\underset{\beta}{\rightleftarrows}}\;
  \operatorname{Hom}(X,Y)\times\operatorname{Hom}(X,Z).
\]
Fie $p_Y:Y\times Z\to Y$ și $p_Z:Y\times Z\to Z$ proiecțiile
canonice:
\[
 p_Y(y,z)=y,\qquad p_Z(y,z)=z.
\]
Diagrama proiecțiilor și a aplicațiilor componente este
\[
\begin{tikzcd}[column sep=large, row sep=large]
 & X \arrow[dl, "h"'] \arrow[dr, "g"]
     \arrow[d, dashed, "{\beta(g,h)}"] & \\
 Z & Y\times Z \arrow[l, "p_Z"] \arrow[r, "p_Y"'] & Y
\end{tikzcd}
\]
Pentru $f:X\to Y\times Z$, definim
\[
  \alpha(f)=(p_Y\circ f,\,p_Z\circ f).
\]
Reciproc, pentru $g:X\to Y$ și $h:X\to Z$, punem
\[
  \beta(g,h)(x)=(g(x),h(x)),\qquad x\in X.
\]

\medskip
\noindent\emph{Observație intercalată în manuscris.} Fixăm $a\in Z$ și
considerăm
\[
 \iota_Y:Y\longrightarrow Y\times Z,\qquad
 \iota_Y(y)=(y,a).
\]

\medskip
\noindent Arătăm mai întâi că $\alpha\circ\beta=\operatorname{id}$.
Pentru orice $g:X\to Y$ și $h:X\to Z$,
\[
  \alpha\bigl(\beta(g,h)\bigr)=(g,h).
\]
Într-adevăr,
\[
\begin{aligned}
 \bigl(p_Y\circ\beta(g,h)\bigr)(x)
   &=p_Y\bigl(\beta(g,h)(x)\bigr) \\
   &=p_Y(g(x),h(x))=g(x),
 \qquad \forall x\in X,\\[3pt]
 \bigl(p_Z\circ\beta(g,h)\bigr)(x)
   &=p_Z\bigl(\beta(g,h)(x)\bigr) \\
   &=p_Z(g(x),h(x))=h(x),
 \qquad \forall x\in X.
\end{aligned}
\]
Prin urmare,
\[
  (\alpha\circ\beta)(g,h)
    =(p_Y\circ\beta(g,h),\,p_Z\circ\beta(g,h))
    =(g,h).
\]

\medskip
\noindent Reciproc, demonstrăm că $\beta\circ\alpha=\operatorname{id}$.
Fie $f:X\to Y\times Z$. Atunci
\[
 (\beta\circ\alpha)(f)=f.
\]
Pentru $x\in X$, notăm $f(x)=(y,z)$. Avem
\[
\begin{aligned}
 \beta\bigl(\alpha(f)\bigr)(x)
   &=\beta(p_Y\circ f,\,p_Z\circ f)(x)\\
   &=\bigl((p_Y\circ f)(x),\,(p_Z\circ f)(x)\bigr)\\
   &=\bigl(p_Y(y,z),\,p_Z(y,z)\bigr)\\
   &=(y,z)=f(x).
\end{aligned}
\]
Deci cele două corespondențe sunt inverse una alteia.
; fără preambul sau pagină nouă.
% NOTE EDITORIALE:
% - Pagina continuă problema 2(b) din Seminarul 3 - Pagina 1.
% - Manuscrisul alternează notațiile P_y/P_z și p_y/p_z;
%   proiecțiile sunt notate consecvent p_Y și p_Z.
% - Diagrama rearanjează schița manuscrisului pentru lizibilitate,
%   păstrând proiecțiile și cele două aplicații componente.
% - O observație intercalată pare să fie: fixez a din Z,
%   i_Y(y)=(y,a). Am păstrat-o, deși nu intervine în demonstrație.
% - Calculele tăiate din dreptul celor două identități sunt omise;
%   pașii finali lizibili sunt păstrați.
% - În câteva locuri manuscrisul scrie identitatea fără a preciza
%   mulțimea pe care acționează; domeniile sunt explicite mai jos.

\subsection*{Seminar 3 - Pagina 2}

\noindent\textbf{(b) Continuare.} Avem corespondența
\[
  \operatorname{Hom}(X,Y\times Z)
  \;\overset{\alpha}{\underset{\beta}{\rightleftarrows}}\;
  \operatorname{Hom}(X,Y)\times\operatorname{Hom}(X,Z).
\]
Fie $p_Y:Y\times Z\to Y$ și $p_Z:Y\times Z\to Z$ proiecțiile
canonice:
\[
 p_Y(y,z)=y,\qquad p_Z(y,z)=z.
\]
Diagrama proiecțiilor și a aplicațiilor componente este
\[
\begin{tikzcd}[column sep=large, row sep=large]
 & X \arrow[dl, "h"'] \arrow[dr, "g"]
     \arrow[d, dashed, "{\beta(g,h)}"] & \\
 Z & Y\times Z \arrow[l, "p_Z"] \arrow[r, "p_Y"'] & Y
\end{tikzcd}
\]
Pentru $f:X\to Y\times Z$, definim
\[
  \alpha(f)=(p_Y\circ f,\,p_Z\circ f).
\]
Reciproc, pentru $g:X\to Y$ și $h:X\to Z$, punem
\[
  \beta(g,h)(x)=(g(x),h(x)),\qquad x\in X.
\]

\medskip
\noindent\emph{Observație intercalată în manuscris.} Fixăm $a\in Z$ și
considerăm
\[
 \iota_Y:Y\longrightarrow Y\times Z,\qquad
 \iota_Y(y)=(y,a).
\]

\medskip
\noindent Arătăm mai întâi că $\alpha\circ\beta=\operatorname{id}$.
Pentru orice $g:X\to Y$ și $h:X\to Z$,
\[
  \alpha\bigl(\beta(g,h)\bigr)=(g,h).
\]
Într-adevăr,
\[
\begin{aligned}
 \bigl(p_Y\circ\beta(g,h)\bigr)(x)
   &=p_Y\bigl(\beta(g,h)(x)\bigr) \\
   &=p_Y(g(x),h(x))=g(x),
 \qquad \forall x\in X,\\[3pt]
 \bigl(p_Z\circ\beta(g,h)\bigr)(x)
   &=p_Z\bigl(\beta(g,h)(x)\bigr) \\
   &=p_Z(g(x),h(x))=h(x),
 \qquad \forall x\in X.
\end{aligned}
\]
Prin urmare,
\[
  (\alpha\circ\beta)(g,h)
    =(p_Y\circ\beta(g,h),\,p_Z\circ\beta(g,h))
    =(g,h).
\]

\medskip
\noindent Reciproc, demonstrăm că $\beta\circ\alpha=\operatorname{id}$.
Fie $f:X\to Y\times Z$. Atunci
\[
 (\beta\circ\alpha)(f)=f.
\]
Pentru $x\in X$, notăm $f(x)=(y,z)$. Avem
\[
\begin{aligned}
 \beta\bigl(\alpha(f)\bigr)(x)
   &=\beta(p_Y\circ f,\,p_Z\circ f)(x)\\
   &=\bigl((p_Y\circ f)(x),\,(p_Z\circ f)(x)\bigr)\\
   &=\bigl(p_Y(y,z),\,p_Z(y,z)\bigr)\\
   &=(y,z)=f(x).
\end{aligned}
\]
Deci cele două corespondențe sunt inverse una alteia.
; fără preambul sau pagină nouă.
% NOTE EDITORIALE:
% - Pagina continuă problema 2(b) din Seminarul 3 - Pagina 1.
% - Manuscrisul alternează notațiile P_y/P_z și p_y/p_z;
%   proiecțiile sunt notate consecvent p_Y și p_Z.
% - Diagrama rearanjează schița manuscrisului pentru lizibilitate,
%   păstrând proiecțiile și cele două aplicații componente.
% - O observație intercalată pare să fie: fixez a din Z,
%   i_Y(y)=(y,a). Am păstrat-o, deși nu intervine în demonstrație.
% - Calculele tăiate din dreptul celor două identități sunt omise;
%   pașii finali lizibili sunt păstrați.
% - În câteva locuri manuscrisul scrie identitatea fără a preciza
%   mulțimea pe care acționează; domeniile sunt explicite mai jos.

\subsection*{Seminar 3 - Pagina 2}

\noindent\textbf{(b) Continuare.} Avem corespondența
\[
  \operatorname{Hom}(X,Y\times Z)
  \;\overset{\alpha}{\underset{\beta}{\rightleftarrows}}\;
  \operatorname{Hom}(X,Y)\times\operatorname{Hom}(X,Z).
\]
Fie $p_Y:Y\times Z\to Y$ și $p_Z:Y\times Z\to Z$ proiecțiile
canonice:
\[
 p_Y(y,z)=y,\qquad p_Z(y,z)=z.
\]
Diagrama proiecțiilor și a aplicațiilor componente este
\[
\begin{tikzcd}[column sep=large, row sep=large]
 & X \arrow[dl, "h"'] \arrow[dr, "g"]
     \arrow[d, dashed, "{\beta(g,h)}"] & \\
 Z & Y\times Z \arrow[l, "p_Z"] \arrow[r, "p_Y"'] & Y
\end{tikzcd}
\]
Pentru $f:X\to Y\times Z$, definim
\[
  \alpha(f)=(p_Y\circ f,\,p_Z\circ f).
\]
Reciproc, pentru $g:X\to Y$ și $h:X\to Z$, punem
\[
  \beta(g,h)(x)=(g(x),h(x)),\qquad x\in X.
\]

\medskip
\noindent\emph{Observație intercalată în manuscris.} Fixăm $a\in Z$ și
considerăm
\[
 \iota_Y:Y\longrightarrow Y\times Z,\qquad
 \iota_Y(y)=(y,a).
\]

\medskip
\noindent Arătăm mai întâi că $\alpha\circ\beta=\operatorname{id}$.
Pentru orice $g:X\to Y$ și $h:X\to Z$,
\[
  \alpha\bigl(\beta(g,h)\bigr)=(g,h).
\]
Într-adevăr,
\[
\begin{aligned}
 \bigl(p_Y\circ\beta(g,h)\bigr)(x)
   &=p_Y\bigl(\beta(g,h)(x)\bigr) \\
   &=p_Y(g(x),h(x))=g(x),
 \qquad \forall x\in X,\\[3pt]
 \bigl(p_Z\circ\beta(g,h)\bigr)(x)
   &=p_Z\bigl(\beta(g,h)(x)\bigr) \\
   &=p_Z(g(x),h(x))=h(x),
 \qquad \forall x\in X.
\end{aligned}
\]
Prin urmare,
\[
  (\alpha\circ\beta)(g,h)
    =(p_Y\circ\beta(g,h),\,p_Z\circ\beta(g,h))
    =(g,h).
\]

\medskip
\noindent Reciproc, demonstrăm că $\beta\circ\alpha=\operatorname{id}$.
Fie $f:X\to Y\times Z$. Atunci
\[
 (\beta\circ\alpha)(f)=f.
\]
Pentru $x\in X$, notăm $f(x)=(y,z)$. Avem
\[
\begin{aligned}
 \beta\bigl(\alpha(f)\bigr)(x)
   &=\beta(p_Y\circ f,\,p_Z\circ f)(x)\\
   &=\bigl((p_Y\circ f)(x),\,(p_Z\circ f)(x)\bigr)\\
   &=\bigl(p_Y(y,z),\,p_Z(y,z)\bigr)\\
   &=(y,z)=f(x).
\end{aligned}
\]
Deci cele două corespondențe sunt inverse una alteia.
